{\LARGE\bfseries\sffamily\color{inkblue} 微重力血管内三维磁驱动柔性机器人\par 统一研究框架\par}
\vspace{.5em}
{\large 物理建模、复值学习、Birkhoff 离散与条件性 KAM 分析\par}
{\small 可扩写论文源码；以第二阶段审查及其冻结证据为依据，2026 年 10 月 29 日。}

\section{研究范围、模型层次与主张状态}
\textbf{用途。} 本框架组织 1025 的研究主线，保留三维血管内磁机器人的最终目标。它是可继续撰写的论文骨架，明确区分已推导恒等式、条件性数学结论、有限数值观察和待完成证明；本文件没有新增实验。中英稿调用同一符号与公式源，新公式标签不替代原稿编号。

\begin{tabularx}{\textwidth}{@{}P{.2\textwidth}XX@{}}
\toprule
层次 & 对象与保留物理 & 当前证据边界\\
\midrule
完整研究目标 & 三维 Cosserat 杆、磁场／电路、脉动血流、柔性血管壁、接触摩擦及微重力边界。 & 已给方程和能量接口；双向流固耦合、生理与硬件标定仍待完成。\\
高维验证层 & 16／32／64 段空间杆；64 段保留姿态、伸长剪切、双向弯曲与扭转。 & 已有规定流场、刚性管腔接触的有限仿真；无连续体或临床证书。\\
理论与学习层 & 12 位形自由度、24 相维数的三维 Kirchhoff 应变 Ritz 族；含变质量和磁端口。 & 已有可微几何模型与训练网络；ROM 误差及一般 Birkhoff 图未认证。\\
\bottomrule
\end{tabularx}

旧两模态 v10 保留为回归基准及其自身冻结证书的对象。新增 8／16 自由度数学试验检验具有明示坐标构造的算法，不能代替 12 自由度几何 ROM 或 64 段空间杆验证。

\subsection{共享符号与主张清单}
\begin{tabularx}{\textwidth}{@{}P{.22\textwidth}X@{}}
\toprule
\lang{Symbol}{符号} & \lang{Meaning and convention}{含义与约定}\\
\midrule
$\bm r\in\R^3,\ \Rot\in\SO(3)$ & \lang{Centerline; material-to-spatial orientation.}{中心线；由材料坐标到空间坐标的姿态。}\\
$\nu,\kappa,\Omega$ & \lang{Shear--extension, curvature--twist, body angular velocity.}{剪切伸长应变、曲率扭率、体角速度。}\\
$(q,p),\ H$ & \lang{Mechanical coordinates, canonical momenta, real energy.}{机械位形、正则动量、实机械能。}\\
$(Q,P),\ h,\tau$ & \lang{Canonically scaled variables, energy and time; $z=(Q,P)$.}{带共同作用量尺度的正则变量、能量与时间；$z=(Q,P)$。}\\
$x=Dz,\ \xi$ & \lang{Statistical input normalization; complex input. Neither changes the physical state dimension.}{统计归一化输入与复输入；不改变物理状态维数。}\\
$\phi(q),\Gmag(q),u$ & \lang{Magnetic potential features, force port and realizable field/coil input.}{磁势特征、广义力端口及可实现的场／线圈输入。}\\
$\Pfaff,\Birk,\Kform$ & \lang{Pfaff coefficient, Birkhoff scalar, closed nondegenerate two-form matrix.}{Pfaff 系数、Birkhoff 标量、闭且非退化二形式矩阵。}\\
$\chi_t,\ S_h$ & \lang{Compatible chart and central generating function.}{相容坐标图与中央生成函数。}\\
$E,S,e,\Mtherm$ & \lang{Total energy, entropy, internal energy, symmetric dissipation operator.}{总能量、熵、内能及对称耗散算子。}\\
\bottomrule
\end{tabularx}



\textbf{状态标签。} “已推导”指在明确模型假设下得到的恒等式或论证；“条件成立”指已有数学结果但机器人应用尚须核验假设；“有限观察”指已有计算记录；“待证”指尚未完成的证明或验证。\path{claims.tsv} 逐项记录路线、假设、证据与下一项证明责任。

\PaperPage
\section{三维杆与统一磁能接口}
以材料坐标 $s\in[0,L]$ 标记截面，线载荷和惯量均采用每单位参考长度测度。空间杆保留两个弯曲方向和扭转；只有独立的 Kirchhoff 分析层才约束伸长与剪切 \citep{gazzola2018}。
\EqKinematics
\EqRodBalance
式\eqref{eq:rod-balance} 假定参考截面角惯量 $\Inertia$ 不随形变和时间变化；否则必须使用完整角动量导数。由同一个本构能量导出内力和内矩：
\EqConstitutive
圆形各向同性截面可取 $C_\nu=\diag(k_sGA,k_sGA,EA)$、$C_\kappa=\diag(EI,EI,GJ_t)$，其中 $J_t=2I$。Kirchhoff 约束为 $\nu=e_3$，内力成为约束反力。夹持运动基座须同时规定位置和姿态；端部条件从边界虚功得到。已运行基准采用固定基座和分布式远端磁段。

\subsection{同一磁势固定力、矩与功率符号}
对于体磁矩固定的硬磁材料，以及截面尺度上变化足够缓慢的外场，
\EqMagneticEnergy
固定姿态对位置变分得到磁力；由 $\delta\mu=\delta\vartheta\times\mu$ 得到磁矩。同一远端磁段不能再重复记作另一个点偶极。Maxwell 表面应力是更一般的总力描述，不可无说明地换成线积分。诱导磁化、磁滞及退磁场需要相应自由能。\citet{tong2025} 的实例提供选定几何与磁性尺度，不提供空间站生理参数。
\EqMagneticPort
局部无电流场具有三个场分量和五个对称无迹梯度分量，但它们不是八个任意模态广义力。线圈场映射、电流与电流变化率约束进一步限制 $u$。已有 12 线圈 Biot--Savart 接口是工程模型，不是实物测量标定。
\Status{已推导；P02--P03}{几何质量和磁虚功来自对应能量；其材料参数、磁化假设和外场适用尺度决定结论范围。}

\PaperPage
\section{血流、接触、微重力与几何降阶}
\subsection{完整目标及已经运行的闭合层}
完整目标包括移动流域中的不可压缩流动、界面速度连续及杆／壁牵引反作用：
\EqFluid
这里的时间导数位于 ALE 坐标。若已求解流固耦合，界面牵引已经给出流体载荷，不能重复叠加 Morison 拖曳。当前仿真采用规定曲管脉动剖面与相对速度阻力，是单向闭合层，并非已求解的无散流场；相关运行关闭了附加质量。\citet{womersley1955} 的直管分析可参照脉动相位，不能证明当前弯管近似。

用有符号间隙 $g$ 定义活动接触区的代表性罚能和切向正则化摩擦：
\EqContact
零切向速度时最后一项定义为零。接触点速度包含平动和转动贡献，切向力同时产生力矩。该指数是 Hertz 型，但采用的线罚系数并未宣称由薄血管壁的 Hertz 解推出。接触开闭非解析，不能直接纳入经典解析 KAM 定理。

微重力既通过残余／平台加速度进入，也通过入口脉动、出口阻抗、顺应性和壁参数等生理边界进入。把重力常数置零不能确定这些数据。规定流体和运动管壁可以对机器人做正功，不能一概替换为 $-C\dot q$ 并宣称机械能单调下降。

\subsection{几何 ROM 与变质量 Hamiltonian}
构形嵌入为 $(\bm r(q,s),\Rot(q,s))$；$J_x,J_\Omega$ 分别将广义速度映射为平动速度和体角速度，因此
\EqMass
对于固定基座嵌入，Legendre 变换给出
\EqHamiltonian
质量导数项由 $\partial_iM^{-1}=-M^{-1}(\partial_iM)M^{-1}$ 得到，不能省略。运动基座会增加速度仿射项和显式时间功，扩展时必须保留。各层使用一致的功率核算：
\EqPower
此式的机械能不含磁势，右侧计入机械磁功；若改把磁势计入能量，则右侧须改用相应的外场显式变化功，不能重复计数。

\PaperPage
\section{方向 A：相容 Birkhoff 表示与生成图}
\subsection{局部存在的条件与实际计算责任}
先固定真实偶维光滑场 $\dot z=f(z,t)$，包含声明的控制和耗散。为区分姿态矩阵，Pfaff 系数记为 $\Pfaff$，不再记作 $R$：
\EqBirkhoff
作用量变分得到最后的动力学关系。反对称及可逆还不够，空间二形式必须闭合，并满足时间输运相容关系：
\EqCompatibility
一种局部构造是求解一形式输运：
\EqTransport
因为 $\nabla(\Pfaff\T f)=(D\Pfaff)\T f+(Df)\T\Pfaff$，直接代入可得式\eqref{eq:birkhoff}。光滑局部微分同胚流保留初始非退化性。这是有条件的存在论证；求输运、图及其导数仍有实际成本。规范变换 $\Pfaff\mapsto\Pfaff+\nabla\sigma$、$\Birk\mapsto\Birk-\sigma_t$ 不能消除原场表示残差。

\subsection{原稿 Theorem A／B 可继承的准确衔接}
相应依据是 \citet{fengqin2010} 第 15 章 Theorem 3.5、4.1、4.4、4.5。端点变换 $\alpha:(z_+,z_-)\mapsto(\widehat w,w)$ 必须适配两端的二形式；令其导数有四个分块，则
\EqEndpoint
$\mathbb J$ 是与 $(\widehat w,w)$ 排列相应的常反对称矩阵。演化中要求 $\det(C_\alpha F+D_\alpha)\ne0$，还原端点也须有局部可逆分支。初始对角线上有
\EqInitialGenerating
初始生成梯度由 $\alpha$ 决定，不能一律设为零。局部存在 $\widehat w=\nabla_w\varphi$，不要求 $N^{-1}$。固定 $w$ 和初始端点得到
\EqGeneratingTime
解析系数递推必须维持梯度性，并通过截断／还原定理得到端点格式。对截断梯度 $G_h(w)=\sum_{j=0}^{m}h^jf_j(w)$，端点还原还要求 $\TruncatedTransversality$，其中 $N_h=D_wG_h$。时变图中的场为 $\partial_t\chi_t+D\chi_t f$，遗漏图的时间导数会改变动力学。
\Status{条件成立 A01--A03；机器人接入待证}{一般受控图推进器尚未接到真实几何 12 自由度 ROM。64 段仿真或已知线性共轭例不能替代这一缺口。}

\PaperPage
\section{方向 A：原（36）、阶数与可实现控制}
一般变质量模型的正确中点方程保留全部质量导数，并同时求解中点未知量：
\EqMidpoint
当 $M,K_s,C$ 为常矩阵、$V=q^TK_sq/2$ 且 $Q_{\rm nc}=F-CM^{-1}p$ 时，消元得到\Original{36}的矩阵形式：
\EqOriginalThirtySix
逆矩阵从左乘，矩阵因子一般不能交换。这是同一个中点方法的重写，不是另一种不同阶数的算法。状态相关系数仍隐含未知中点；实际应解线性系统或 Newton 残差，不必显式求逆。

在光滑稳定分支上中点为二阶。自治正则 Hamiltonian 的中央生成函数修正为
\EqCentral
恒等式 $J\nabla(f^TH''f)=D^2f[f,f]-2(Df)^2f$ 给出三次增量及四阶。这里的正号用于预修正，不能与普通中点后向修正 Hamiltonian 中的负号混淆。所需假设包括足够光滑、导数对称和局部唯一隐式分支；有限 Newton 残差另计实施误差。一自由度阻尼振子给出明确反例：
\EqDampingCounterexample
因此普通阻尼中点既非通常正则辛格式，也不自动保持所选指数权重。分块加权图要求相容的等阻尼率分块及势能，任意跨块耦合会破坏其特殊构造。

\subsection{磁实现是独立于算法的控制条件}
动量方程中的目标力 $U$ 以 $hU$ 的尺度进入时间步；它必须通过实际磁端口分配：
\EqAllocation
一类可积反馈具有一致的闭环势：
\EqPotentialControl
真实磁力对位形求偏导时固定输入。直接把反馈代入原磁势再对整个复合函数求导，一般会增加非物理链式项；式\eqref{eq:potential-control} 是另外构造的正确闭环势。

\PaperPage
\section{方向 A：结构、表示和收敛分别验收}
新增场 $g$ 保持同一闭二形式的一个局部充分条件，以及采样反馈实际需要的导数，为
\EqControlDerivative
势形反馈须按闭环势一致取梯度并离散。固定电流子步可以是辛的，整个状态反馈映射却未必辛。共点阻尼 $y=G_m^T\dot q$、$u=-K_yy$ 给出非正控制功率；这是被动性，不能称为通常辛性。

\subsection{标明残差对象的条件性估计}
近似一形式准确表示的场可能不同于目标场：
\EqRepresentationError
假设轨道位于共同紧域，原场光滑且 Lipschitz，隐式分支稳定、图及其逆有界，$hL_f/2<1$，状态增量残差 $\|r_k\|\le\eta_h$，且有统一界 $\|\Kform_a^{-1}\|\le\kappa_B$、$\|e_B\|\le\varepsilon_B$。离散稳定估计及 Gronwall 给出
\EqConvergence
第一行属于中点。第二行取 $p=2$；只有另证五阶局部余项和阶段稳定性后才可取 $p=4$。阶段残差须先换算为状态残差。接触切换、限幅和零阶保持均可能改变阶数。

若近似图具有下述缺陷，而内部理想映射精确保持 $K_0$，外部拉回映射的结构缺陷为
\EqChartError
内部映射仅近似求解时，还需叠加实施与切线计算误差。这个恒等式和轨道误差界回答不同的问题。

\subsection{已有数值证据的确切对象}
8／16 自由度算法试验采用明示的已知正则点提升，并非任意杆质量度量。$n=16,T=2,h=0.025$ 时，保守中点与 Central4 的观测阶为 1.9992、3.9994；分块加权阻尼下，合法加权中点与非自治图 Gauss4 为 1.9992、3.9998。一般 ODE Gauss4 不能因此在耗散场上获得精确保结构结论。

一形式／Moser 试验在一个 16 自由度状态、$t=0.1$ 处，图匹配缺陷为 $6.68\times10^{-16}$，但原场误差为 $0.109619$。极小结构残差没有认证原动力学准确。另一组独立诊断所用的精确流信息未进入状态推进；已知共轭和辅助图计算成本均有记录。
\Status{条件性推导 A06--A07；有限观察 A08}{证据见 \path{../notes/route_a.md}；尚无完整机器人工作域的统一界或新 KAM 证书。}

\PaperPage
\section{方向 B：机械学习、正规形与条件性 KAM}
12 自由度 Kirchhoff Ritz 模型保留三维弯曲和扭转，并由几何得到 $M(q)$。正则尺度化固定 $h_2$，其中 $\nu_j=\omega_j/\omega_0$；统计归一化是另一项操作：
\EqLearningNormalization
完整复网络采用可训练复隐藏层、共轭输入、实能量输出、实／虚部分别取 $\tanh$、三次零偏置门控及动量偶对称。小经验梯度误差不是连接学习与 KAM 所需的复解析域界 \citep{chen2022}。原非受限网络已出现负 $h_{PP}$ 样本，因而对应采样域上的全域正质量断言已有反例。

另行训练的机械架构将动量依赖限制为正定二次型，不减少状态维数：
\EqMechanicalNetwork
质量网络减去原点输出使 $C(0)=I$；三次势门控固定残差 0／1／2 阶 jet。对非零 $v$，有 $v^Th_{PP}v=\|L^TD_\nu v\|^2>0$，证明了逐点动量可逆性，但没有给出可用的统一条件数和逼近界。

\begin{tabularx}{\textwidth}{@{}Xrr@{}}
\toprule
新确认集，种子 1025 & 原非受限 $\nabla_z h$ 误差 & SPD $\nabla_z h$ 误差\\
\midrule
主域 test／OOD & 1.950\%／4.693\% & 6.856\%／13.608\%\\
压力域 test／OOD & 6.043\%／9.463\% & 98.883\%／119.752\%\\
\bottomrule
\end{tabularx}
SPD 网络有 13,440 个实参数，原非受限网络为 7,104，不能作等容量表示比较。压力域样本条件数约达 $10^{12}$，该不利结果完整保留。已有前馈使用已知几何 $M$，仅作受限的学习 $h_Q$ 补偿，没有反解失效的神经质量。其结构阻尼系数明确为 $\alpha_s=2\,\mathrm{s}^{-1}$，时间尺度化后为 $\alpha_s/\omega_0$。

\subsection{连续和离散 KAM 的待证责任不同}
圆形各向同性杆与轴向偏场产生四组 1:1 弯曲频率，不能照搬非共振同调除法。先研究共振正规形，或基于真实物理证据建立破对称模型。后验定理还需指定环面、解析域、Diophantine 频率、非退化性及完整常数 \citep{haroluque2019}；一个已存在非退化环面附近的稳定性不能分类任意工作域 \citep{bounemoura2016}。
\EqKAMDivisors
第二行对应 \citet{shang1999} 所选小扭转离散条件，不是替代该定理全部假设的通用式；它对所有非零整数 $k$ 成立，并需频域边界余量、扭转与小量阈值。连续非共振不自动推出固定步长上的离散条件。

\textbf{先固定精确映射。} 离散 KAM 研究一个明确的理想解析 exact-symplectic 映射，还需统一解析隐式分支和相容图。求解／图误差只有在另证受扰映射仍属于这一类时，才可计入生成扰动预算。任意小的非辛实施残差、采样结构缺陷或 RMSE 都不能直接相加成为 Hamiltonian 的 KAM 预算。原 v10 证书只属于原冻结两模态保守 Hamiltonian。

\PaperPage
\section{Metriplectic 候选与后续证明责任}
机械 Rayleigh 损耗和外部功率端口足以定义已经运行的机械试验；只有把热效应本身列为目标时，热状态才成为必要建模对象。Metriplectic 将可逆流、对称不可逆交换和退化条件分开 \citep{coquinot2020}，不能把它视为 Birkhoff 几何的另一种名称。

\subsection{已推导的有限维候选}
对于自治机械能及内部热库，取
\EqThermalState
\EqMetriplecticMatrices
对称因式分解给出正半定性，直接相乘给出两个退化恒等式：
\EqDegeneracy
闭系统相应满足
\EqThermalBalance
机械能经 $-Cv$ 损耗后进入内能。开放系统的磁场、流体、管壁和基座功率仍须另加。改用熵坐标 $e=U(S)$ 时，算子按坐标变换处理；温度因子的外观变化不是模型矛盾。奇维 Poisson 矩阵本来退化，不能求逆成为全状态 Birkhoff 二形式。

该候选已经在真实 12 自由度 ROM 的 32 个状态上作代数核验，但不是温升模拟：热容量、传导、血流换热及线圈／磁滞损耗均未识别。规定运动流体的机械功率必须单独核算，不能直接全部变成机器人内热源 $v^TCv$。

\subsection{主张清单与扩写顺序}
\begin{tabularx}{\textwidth}{@{}P{.14\textwidth}P{.23\textwidth}X@{}}
\toprule
路线 & 范围内已有内容 & 更强论文结论的前置任务\\
\midrule
A & 条件性恒等式、合法特殊算法、8／16 自由度有限数据。 & 在几何 ROM 上构造一般图；统一表示／图／求解界；真实磁分配与事件处理。\\
B & 实能量复学习、固定 jet、SPD 构造及不利拟合结果。 & 准确且条件良好的机械类；共振正规形；独立同域界及新环面证书。\\
物理 & 64 段空间杆验证及共享能量／磁接口。 & 线圈和生理标定、流体／壁耦合、截断及空间时间误差控制。\\
热力学 & 正半定算子候选和退化证明。 & 明确热研究目标并补齐热库参数，再进行完整热部署。\\
\bottomrule
\end{tabularx}

\Next{先验收模型与执行器工作域，再在该域辨识 ROM 与学习误差。方向 A 先分别测量表示与图误差，再作长期推进；方向 B 先冻结保守解析候选及选定定理，再计算常数。接触与血流始终保留在最终目标中；光滑无接触邻域只是有条件的分析分支。}

\textbf{来源责任。} 物理、学习、Birkhoff 与综合报告，以及 \path{../notes/route_a.md}、\path{../notes/learning_spd.md} 提供现有证据。\path{claims.tsv} 将主张逐项连接到记录；\path{reference_verification.tsv} 保存 DOI 核验页面。引用文献不等于其定理已用于新机器人。

\PaperPage
\small
\bibliographystyle{plainnat}
\bibliography{references}
